% =====================================================================
%  締め：実務への手順，考察，限界，まとめ，投稿までに残っている作業
%  文の元：paper/sections/7-discussion.tex，8-limitations-ethics.tex，9-conclusion.tex
%  残っている作業：2026-10-27 の組版（本文 12 頁）と，この資料を作るときに見つけた論文の誤り
% =====================================================================
\setband{締め}

\begin{frame}{新しいコーパスに数本のラベルしかないときの 4 つの手順}
  \rl{考察：実務への手順}
  \vspace{0.4em}
  \centering
  \begin{tikzpicture}[font=\scriptsize]
    \foreach \i/\t/\d/\c in {
      0/中心化/{ソースとターゲットを別々に中心化する．ラベル無しで判断基準の損の大半が消える}/cCrit,
      1/全ソースから出発/{全ソースからの出発点に，切片はゆるく，重みは強く縛る}/cStart,
      2/選び方に頼らない/{ラベルを付ける発話をラベル無しで選んでも，よい出発点の向きは直らない}/mMuted,
      3/別のモダリティ/{別のモダリティが記録されていれば（ラベルを付けるときだけでも），ソースで学んだ別のモダリティの先生にプール全体のラベルを付けさせる}/cExt} {
      \node[draw=\c, line width=0.9pt, fill=white, rounded corners=3pt, text width=3.4cm, minimum height=3.0cm, align=left, anchor=north] (b\i) at (\i*3.85,0) {};
      \node[anchor=north west, font=\footnotesize\bfseries, text=\c] at ($(b\i.north west)+(0.12,-0.1)$) {\number\numexpr\i+1\relax．\t};
      \node[anchor=north west, text width=3.15cm, align=left] at ($(b\i.north west)+(0.12,-0.6)$) {\d};
    }
    \foreach \i in {0,1,2} \draw[arr] (b\i.east) -- ++(0.4,0);
  \end{tikzpicture}\par
  \vspace{0.6em}\raggedright
  \begin{itemize}\scriptsize\setlength{\itemsep}{4pt}
        \item 4 でも，テストのときは生徒のモダリティだけを使う
      \end{itemize}
\end{frame}

\begin{frame}{考察：言葉の振る舞いと，クラスタの仮定に頼らない弱いラベル}
  \rl{考察}
  \vspace{0.3em}
  \begin{itemize}\footnotesize\setlength{\itemsep}{10pt}
        \item \textbf{言葉}：判断基準の損が小さいので，ラベルが直す判断基準もほとんど無い（割合は最大 \NUM{p1-e29-critshare-k1-t}）．それでも融合の先生から向きの情報を得て上がる
        \item \textbf{エントロピー最小化}：感情がラベル無しのデータで塊を作るという仮定に頼る．この仮定は IEMOCAP 即興の AVT・AT でしか成り立たない
        \item \textbf{モダリティをまたぐ弱いラベル}：モダリティの相補性に頼る．\NUM{st-b-tgt-n} の組み合わせのうち，拡張版が出発点を下回るのは \NUM{st-b-tgtneg-n}（最悪 \NUM{st-b-tgtneg-worst} \pt）
      \end{itemize}
  \hosoku{s:c:disc}{考察の細部}
\end{frame}

\begin{frame}{限界：凍結した特徴と線形分類器，近似に立つ式，4 感情と 6 つの英語のコーパス}
  \rl{限界}
  \vspace{-0.2em}
  \begin{itemize}\scriptsize\setlength{\itemsep}{4pt}
        \item \textbf{範囲}：凍結したエンコーダと線形分類器だけ．微調整の確認は音声 1 つ，4 ターゲット，反復 3 回
        \item \textbf{式}：2 つの山の模型（正規分布，共通の分散）と 1 次・2 次の近似に立つ．予言を実測で確かめたが，厳密な結果とは言わない
        \item \textbf{判断基準が先}：単独の音声と顔の $K\le5$ に限る．顔で割合が大きい理由は説明できていない
        \item \textbf{データ}：4 感情，英語の 6 コーパス．CMU-MOSEI の話者の分離は不完全．言葉を含む比較は 4 ターゲット
        \item \textbf{統計と選び方}：区間はコーパスの中の抜き方のばらつきだけ．仮のターゲットで選ぶと最大 \NUM{cmp-oursx-oracle-vs-tuned-avt-t-k0} \pt 失う
        \item \textbf{列の割り直し}：利得が感情の比にどう依存するかは調べていない（偏りが強いと害になりうる）
        \item \textbf{倫理}：公開のコーパスだけを使った．個人についての判断には検証していない
      \end{itemize}
\end{frame}

\begin{frame}{まとめ：判断基準は少ないラベルで直り，向きは別のモダリティの弱いラベルで直る}
  \rl{まとめ}
  \vspace{0.3em}
  \begin{itemize}\footnotesize\setlength{\itemsep}{10pt}
        \item 少ないラベルはまず判断基準を直す．音声の向きは $q/K$ でしか良くならない
        \item 音声の 1 ソースは，向きについて感情あたり \NUM{p1-e32-n1src-audio-lo}〜\NUM{p1-e32-n1src-audio-hi} 本ぶんを補う．ラベル無しの選び方では向きは動かない
        \item ソースだけで学んだ別のモダリティの先生の弱いラベルは向きを動かし，拡張版はすべての生徒で最も強い比較手法を上回る
        \item 次の一歩：$q$，$e$，$J$ を仮のターゲットで推定し，ラベル・ソース・モダリティのどれを足すかを，ラベルを付ける前に選ぶ
      \end{itemize}
\end{frame}

\begin{frame}{投稿までに残っている作業}
  \rl{残っている作業}
  \vspace{0.3em}
  \begin{itemize}\footnotesize\setlength{\itemsep}{7pt}
    \item \textbf{著者と略歴}：著者名・所属・略歴がまだ空欄．題は仮題
    \item \textbf{頁}：本文は参考文献 78 本を含めて 12 頁で，12 頁目は右の列の下の約 2 割が空いている．略歴を入れると 12 頁を少し超えるおそれがあり，投稿の目安の 11.5 頁には届いていない
    \item \textbf{論文と台本の食い違い}：この資料を作るときに 8 つ見つけた（次のスライド）．論文と \texttt{CLAIM.md} の文を直す必要がある
    \item \textbf{未検証}：列の割り直しの利得と感情の比の関係（R1）は計算していない．計算には承認が要る
  \end{itemize}
\end{frame}

\begin{frame}{論文の書き方が台本や結果と違う点（この資料を作るときに見つけた）}
  \rl{残っている作業}
  \vspace{-0.3em}
  \centering\tiny
  \renewcommand{\arraystretch}{1.12}
  \begin{tabularx}{\textwidth}{@{}l X X@{}}
    \toprule
    論文の箇所 & 論文の書き方 & 台本と結果ファイル \\
    \midrule
    5・6 節，\texttt{CLAIM.md} & 基本版の先生は融合の出発点（融合の生徒には，それが教えた単独の生徒の積） & 単独の生徒の候補に，単独の出発点の予測の積と，一致で重みづけた融合の出発点もある．融合の生徒の候補には，単独の出発点の積と，教えた生徒の積をそのまま予測に使う版もある．どれもよく選ばれる（\texttt{exp/multimodal\_distill.py}，補足 2-12） \\
    5・6 節 & 拡張版は基本版の候補に先生を足す & 候補は $K$ 本で学び直した融合の生徒と，それが教えた生徒の積だけで，基本版の候補を含まない（\texttt{exp/baselines\_distill.py}） \\
    5 節 & 5 つの値は 3 つずつの候補から選ぶ & 第 II 部の後半の A・A(HuBERT)・V・AV の生徒は，E40 で選び直した値（$C'=0.003$ など，3 つの候補の外）を使う（\texttt{exp/geom\_e0.py}） \\
    6 節 & 固定版の EntM・CoR・PA は単独の生徒を 1 つの感情に潰す & UAR が約 25\% まで下がるのは CoR と EntM の音声の生徒だけ．PA は 31〜35\%（補足 3-15） \\
    3 節 & $\Delta$ はターゲットの全ラベルの線形判別で求める & 話者を半分に分け，片方で縮小推定の線形判別の向きを求め，もう片方で測って入れ替えて平均する（\texttt{exp/explore\_pairs.py}） \\
    3 節，補足資料 & 全件では，$C$ の規則が合併の損を大きく見せる & 大きく見せるのは主に音声で，ほかの種類では差が小さい（補足 0-6） \\
    5 節 & the face start falls \NUM{p2-e60pre-juar-av-a-v} points below it & 値が負なので符号が二重になる \\
    3 節 & $o = (\mu_k + \mu_l)/(2\sigma)$ & 台本の符号は逆．$a(d,o)$ は左右対称なので結果は変わらない \\
    \bottomrule
  \end{tabularx}
\end{frame}

